\section{Discussion}\label{sec:discussion}

\subsection{Two weights on growth}
The model separates two weights on growth that practice treats as one. The screening weight of \autoref{lem:valuation2} compares a persistent change in growth with a persistent change in margin; it is a duration, about 1.8 to 2.8 for a firm with a ten-year runway and partly transitory margins, close to the weights of two to three that \citet{bessemer2024rulex} recommend for free-cash-flow margins ``if growth is efficient'', a conservative base within a reported range of about 2 to 9. Under the verification hypothesis (H) of \autoref{prop:dynamic}, the decision weight is the marginal value of a dollar of revenue, the marginal $q$ of revenue in the sense of \citet{hayashi1982tobin}. In the worked example it is close to the revenue multiple, 5 to 15 at a discount rate of 10\%; that is an illustration, not an estimate from the data, and an owner who discounts at a higher rate acts on a lower weight. In the example of Section~\ref{supp:model-weights} of the Supplementary Appendix the weight is \WeightTen{} at 10\% and falls to \WeightTwentyTwo{} at 22\% and \WeightThirtyOne{} at 31\%, within or below the screening range. It exceeds the screening weight when the firm's revenue multiple is well above it, which is a statement about valuation and about the owner's discount rate, not about the frontier. The data do not settle it: the ratio of growth spend to excess growth bounds the net weight from below for a given measure of spend, but the measure moves the ratio from \AlphaAllLo{} to \AlphaAllHi{} (Section~\ref{sec:empirics}). Because cash flows of longer duration earn lower expected returns \citep{weber2018cash,gormsen2023duration}, a constant market discount rate, if anything, understates the decision weight of an owner who discounts at the market's rates. The screening weight tells an investor how to compare firms; under (H), the decision weight tells a manager whether a dollar of growth spend creates value.

The difference has a practical consequence. A firm that acts on its decision weight $X$ spends $\sigma^*=\alpha(X-k)\Delta^*$ (\autoref{prop:optimum}), so its Rule-of-40 score exceeds the saturated level $m_0+(1-k)G$ by $\Delta^*(1-k-\alpha(X-k))$, which is negative whenever $\alpha(X-k)>1-k$. The sign depends on the product $\alpha(X-k)$, which the data do not identify. If all of a firm's selling and marketing expense buys growth and the median spend ratio of \SpendEone{} is its $\alpha(X-k)$, an asset-light firm growing 10 points faster than its market scores about \SpendPenaltyTen{} points below the score it would earn with no growth spend, even though the spending creates value at its own weight; with maintenance spend netted out, a ratio of \AlphaNet, the shortfall is about zero. Among the best firm-years of 2023--25, a quarter score at least \AnchorPiLoAbs{} points below their durable level if all of their selling and marketing expense buys growth (Section~\ref{sec:horizons-ranges}). A fixed threshold on the Rule of 40, whether 40, 50 or 60, therefore marks down such firms when their marginal growth spend earns more than it costs; whether a typical firm's does depends on how much of its selling expense buys growth, which the data do not settle. A firm that claims it must show it from the payback and retention of its customer cohorts at a stated discount rate. \autoref{prop:general} gives the same point in value terms: the Rule of $X$ at the firm's own weight is a criterion of economic profit plus growth at its net value, while the cash Rule of 40 weights growth at $1-k$ whatever it earns.

\subsection{Execution sets the pace}
The static frontier says where a firm can be once a transformation is complete; the attainable frontier says how fast it gets there. In our calibration the binding constraint is the capacity to change the work of employees, not decision bandwidth or security review, and failures, rollbacks, leadership changes and reorganizations slow the program, approximately like a cut in its pace (\autoref{prop:expected}). The difference between central and fast execution, about \PaceCenSix\% of the program's benefits realized by 2036 against an upper bound on the embedded share that reaches the whole program in \PaceFastDone, lies in the share of transitions that create value and the stability of their sponsorship, not in technology, given the assumption that every process is technically feasible from the start. Two implications follow. A firm in the middle of a program scores below its eventual durable score by the dip of the transitions in flight, about \CenDipEight{} to \PaceDipFast{} points in our scenarios, so a screen that reads the dip as weakness penalizes firms that are executing. The program repays the dip only after \PayBackCen{} years under central execution, so whether it is worth starting depends on the owner's horizon only if the owner cannot recover the remaining benefits in a sale or must repay the program in cash within that horizon. And in the cases computed here ownership changes the pace, not the target (Section~\ref{sec:horizons-ranges}). These statements rest on scenario values (Section~\ref{sec:discussion-limits}).

\subsection{Measurement at the limit}
Near the top of the distribution, accounting choices move the score by more than operating performance does. Four mechanisms matter. First, adjusted margins exclude SBC (Palantir's also exclude the related employer payroll taxes): Palantir's second-quarter 2026 score was 155 on its adjusted operating margin of 62\% but 140 on its GAAP operating margin of 47\% \citep{palantir2026q2}. Its annual GAAP score for fiscal 2025 was \PalantirGAAP, but that comparison also reflects an acceleration of growth, not measurement alone. Second, computing capacity that a firm owns is capitalized. It enters operating income only through depreciation over the useful life management assigns, it never enters EBITDA, and it appears in full only in free cash flow. Lengthening the assigned useful life of servers raises operating margins without changing cash flow, and capitalizing internally developed software shifts cost between periods in the same way; for a firm whose capitalized spending grows with revenue, amortization charges today's income with past, smaller vintages, so the GAAP score exceeds the cash score permanently, by more the faster the firm grows \citep{murthy2026transformation}. Third, recording revenue gross rather than net raises reported revenue relative to profit, and it raises reported growth while the mix of gross-recorded revenue is rising. Fourth, run-rate and annualized quarterly revenue run ahead of annual revenue in a growth spurt. Each of these can push a reported score above the full-automation ceiling without any change in economics. The ceiling itself is a bound at a stated cost structure, not on observed scores: a firm whose costs are leaner than that structure, for example with a lower compute share, can score above it in economic terms too (\autoref{cor:ceiling}).

Growth that is bought also needs care: acquired revenue enters a reported score without the capital paid for it, and interest on acquisition debt is not in the operating margin, although under U.S.\ GAAP it is in operating cash flow\IfFileExists{sections/measurement_acquired.tex}{; Section~\ref{sec:acquired-growth} shows how to charge both without charging financing twice}{}. For Bending Spoons, on figures from its F-1 filed on 8 June 2026, which are unchanged in the final prospectus, the 2025 GAAP score of \acqRGaapTwentyFive{} is \acqRIntTwentyFive{} after interest; on the cash margin before interest, with acquisitions charged as capital, the score lies between \acqConvLo{} and \acqConvHi, depending on the charging convention. Such a score is evidence of access to targets and capital, not of a frontier above the ceiling, and a firm whose base shrinks and that wins no new customers must keep buying to stand still. This is why statements about reported scores use observed scores, while a statement about the frontier is made only where it also holds on the organic frontier, which removes acquired growth. We suggest computing the score as this paper does, on annual free cash flow after SBC and after capital expenditure including capitalized software, and, where lease-financed capacity is material, also after additions to finance-lease assets. Scores on this basis are not comparable with the scores practitioners usually quote, so the basis should be stated with every score.

\emph{Use in filings.} A published score on adjusted margins, free cash flow, organic growth or charged acquisitions is a non-GAAP measure, with the reconciliation and prominence that requires, and a Rule-of-$X$ metric in management's discussion and analysis is a key performance indicator; the ceiling, the durable score, the growth surplus, the decision weight and the central path rest on scenario values and are not suitable as reported metrics or targets (Section~\ref{supp:model-filings} of the Supplementary Appendix gives the rules and their sources).

\subsection{What to track if headcount falls}
The Rule of 40 was built for firms whose main costs were people, so its implicit trade-off is between selling and marketing expense and growth. When automation turns labor into capital, capital becomes the scarce input, and the score has no capital term. Two firms can both score 100 while earning very different returns on the capital that produced the score. \autoref{prop:general} says what to track instead: economic profit, defined as after-tax operating profit minus the cost of capital on invested capital including leases and purchase commitments, plus growth at the net weight $X-k$. Two further supplements follow from the model: gross profit per dollar of computing cost, because the compute share is the largest uncertainty in the ceiling; and the growth-weighted score with a weight that reflects the firm's own marginal value of revenue (\autoref{prop:dynamic}, under (H)) rather than a fixed screening weight.

\subsection{The accountability floor}
Full automation does not remove the cost of being accountable. Under Section 404 of the Sarbanes--Oxley Act, management must assess the company's internal control over financial reporting every year, and the auditor of an accelerated or large accelerated filer that is not an emerging growth company must express an opinion on the effectiveness of internal control over financial reporting in an integrated audit performed under PCAOB Auditing Standard 2201, paragraphs .03 and .06 \citep{usc2026sec7262,pcaob2026as2201,cfr2026sec229308}. Emerging growth companies are exempt from the auditor's attestation, and a newly public company's first annual report need not contain management's assessment \citep{usc2026sec7262,cfr2026sec229308}, so the floor of a young public firm is lower than that of a seasoned one. Sections 302 and 906 make the principal executive and financial officers personally responsible for certifying the periodic reports \citep{usc2026sec7241,usc2026sec1350}. Regulators, insurers and courts look to identifiable people and entities. The recurring costs of this system (control testing, the integrated audit, directors' and officers' insurance and the legal and compliance staff behind them) form the floor $a_{\min}$ in \autoref{cor:ceiling}. Much of the floor is fixed in dollars, so as a share of revenue it falls with scale: in the core sample the median general and administrative expense, which contains most of it, is \GAsmall\% of revenue for firms with less than \$1 billion of revenue in 2023--25, \GAmid\% between \$1 and \$10 billion and \GAlarge\% above \$10 billion, and Microsoft and Oracle report \GAleanLo{} to \GAleanHi\%. Section~\ref{sec:horizons-ceiling} sets the band.

Automated controls can lower the cost of testing, but the floor need not fall with automation and may rise with it. Internal control is a pass-fail standard, so \autoref{prop:amin} gives the floor as the cost of the cheapest admissible control design, one that leaves no reasonable possibility of an undetected material misstatement, rather than as a trade of control cost against expected loss. When control layers share one AI model, they share its failures, so the common-mode part of the miss probability does not vanish however many layers are added, and it does not fall with scale. In the illustrative example of \autoref{rem:amin-values} (Section~\ref{sec:horizons-ceiling}), no design built only from AI layers on one provider is admissible at either of the two scenario thresholds. Each AI application that enters the scope of internal control also adds its own IT general controls. Diversity of the control design and the quality of IT general controls, not headcount, therefore decide admissibility, and firms that can bear accountability credibly and cheaply hold a scarce input in the sense of \autoref{prop:entry}. \autoref{prop:amin} is our operationalization of the standard, not an accepted audit method, and the floor omits three things that would raise it. The integrated audit costs the same for every design in \autoref{ass:layers}, although an auditor may benchmark an automated application control when IT general controls are effective and the control is unchanged since its baseline \citep[paragraphs B28 to B31]{pcaob2026as2201}, which a control run by a model that its vendor updates does not meet, so AI-performed controls are likely to need more audit effort. The model prices controls over processes, systems and integrations, but not entity-level and policy controls or the controls over a model provider as a service organization. And it has no parameter for automation bias, the tendency of a reviewer to accept wrong output from an automated aid; across 106 experimental studies, combinations of humans and AI performed worse on average than the better of the two alone \citep{vaccaro2024when}. What remains with accountable people is judgment, which AI's cheaper prediction need not replace \citep{agrawal2019exploring}.

\subsection{Limitations}\label{sec:discussion-limits}
The evidence covers firms that file XBRL financial statements with the SEC under US GAAP, mostly U.S.\ firms. The fastest-growing AI firms of 2025 and 2026 are private and absent, so the sample likely understates the current frontier at high growth; this is one reason the 2028 to 2036 path comes from the calibration rather than from the sample. Industry membership uses each registrant's current SIC code and latest business description; the rule on principal business is a keyword check that misses some game publishers and an AI-cloud provider, with independent readings as a sensitivity, and the rule on listed equity, which could be checked only for 10-K filers, enters only as a sensitivity (Section~\ref{sec:data-rules}). The cross-sectional frontier records the points firms chose, so it does not identify any firm's own trade-off (\autoref{prop:industry}); the data do not identify the spend elasticity $\alpha$ or the weight on growth that firms act on; and the decision weight's interpretation rests on hypothesis (H), which we state but do not prove. The ceiling rests on a compute share that no source measures for a firm's own automation (the nearest figure, total internal AI spending of 11\% of revenue in one survey, is not itemized) and on scenario values for loss severities and the build stock; the pace rests on scenario values for program size, embedding times and the persistence of sponsorship, which no source we found measures; and the central thresholds hold output prices fixed, with pass-through as a scenario. Table~\ref{tab:ai-map} lists channels, some of them scenarios, without evidence on their size, and the model does not net the negative channels for incumbents, such as the loss of seat-priced revenue when customers cut headcount and the entry of AI-native rivals whose growth draws on incumbents' headroom. The calibration's sensitivities move one input at a time; a joint stress test over the stated ranges (Section~\ref{supp:stress} of the Supplementary Appendix) finds that several headline claims hold at every point of the box, that six hold only under their stated conditions, which the text states where each appears, and that the remaining failures lie outside a claim's stated conditions or restate cases the paper reports. The calibration assumes that the work automated by each horizon is technically feasible by then; the paper does not model AI capability, whose measured growth is fast (the length of software tasks that agents complete half the time doubled about every 131 days since 2023; \citealp{metr2026th11}). Finally, we do not test whether any version of the score predicts valuation multiples, returns or default, and we estimate no mapping from scores to multiples; the implications of Section~\ref{sec:conclusion} concern measurement and firm decisions, not investment.
